\documentclass{article}
\usepackage{geometry}
\usepackage{graphicx} % Required for inserting images
\usepackage{amsmath, amsthm, amssymb}
\usepackage{parskip}
\newgeometry{vmargin={15mm}, hmargin={24mm,34mm}}
\theoremstyle{definition} 
\newtheorem{definition}{Definition}

\newtheorem{theorem}{Theorem}[section]
\newtheorem{lemma}[theorem]{Lemma}
\newtheorem{proposition}[theorem]{Proposition}
\newtheorem{corollary}{Corollary}[theorem]


\title{MATH210B Homework 9}
\date{March 2025}
\author{Boran Erol}

\begin{document}

\maketitle

\section{Problem 1}

% \begin{lemma}
%     For any field $ K $ of characteristic $ p $, $ K / K^{p} $ is purely inseparable.
% \end{lemma}
% \begin{proof}
%     Let $ x \in K $ and assume $ x \notin K^{p} $.
%     Notice that $ x $ satisfies $ f(t) = t^{p} - x^{p} = (t - x)^{p} \in K^{p}[x] $,
%     but it doesn't satisfy $ (x - t)^{k} $ for $ k < p $ since otherwise $ x^{k} $ would
%     be in $ K $ which would imply $ x $ is in $ K $, which is a contradiction.
%     Thus, $ x $ is not separable and it's minimal polynomial has degree $ p $ over $ K^{p} $.
% \end{proof}

% Let $ K = F(\alpha_{1}, \ldots, \alpha_{m}) $.

% We'll show that

Let $ \alpha \in K \backslash F $ and consider the minimal polynomial $ p $ of $ \alpha $ over $ F $.

Since $ \alpha \notin F $, $ p $ is not separable. Then, $ p^{\prime} = 0 $.
Then, $ p(x) = f(x^{p}) $ for some $ x \in F[x] $.

Notice that $ f $ also needs to irreducible, so if $ f $ is not separable, we can keep applying this argument to $ f $.

Since $ K/F $ is a finite extension, there's some $ n > 0 $ such that $ p(x) = f(x^{p^{n}}) $ and $ f $ is separable
and irreducible. Then, $ f $ has to split since the extension is purely inseparable. But then, $ f $ has a single linear factor
since otherwise $ f $ would be reducible.

Then, $ f(x) = x^{p^{n}} - a $ for some $ a \in F $. But then $ \alpha^{p^{n}} = a $ which is what we wanted to show.

\section{Problem 2}

Let $ \alpha \in E^{G} $ and assume that $ \alpha $ is separable over $ F $.

Then, $ \alpha \in K $ is also separable over $ F $.

Let $ m_{\alpha} $ be the minimal polynomial of $ \alpha $.

Our goal is to show that $ m_{\alpha} $ is $ x - \alpha $.

$ m_{\alpha} $ splits over $ E $ since it's a normal extension and it already contains
one root, namely, $ \alpha $. Assume by contradiction that it has another root $ \beta $
that is distinct from $ \alpha $.

Then, we have a field isomorphism $ F(\alpha) \to F(\beta) $ by Homework 7 Problem 6.


 
\section{Problem 3}

\begin{lemma}
    Let $ K/F $ be a finite field extension and $ L/F $ be its maximal separable subextension.
    Then, $ K/L $ is purely inseparable.
\end{lemma}
\begin{proof}
    Let $ \alpha \in K $ be separable over $ L $.
    Then, $ m_{\alpha} $ is separable over $ L $, so it's separable over $ F $.
    Then, $ \alpha \in L $ by definition.
\end{proof}

\section{Problem 4}

\begin{lemma}
    Let $ K/F $ be an algebraic field extension. Then, every field homomorphism $ \sigma: K \to K $
    over $ F $ is an isomorphism.
\end{lemma}
\begin{proof}
    $ \sigma $ is injective since $ \sigma $ fixes $ F $ and fields have no non-trivial ideals.

    Let $ \alpha \in K $.

    Let $ m_{\alpha} \in F[x] $ be the minimal polynomial of $ \alpha $.

    Let $ L $ be the subfield of $ K $ generated by the roots of $ m_{\alpha} $ that
    is already in $ K $. Notice that $ L/F $ is a finite extension, and thus the
    restriction of $ \sigma $ to $ L $ is an isomorphism using Rank-Nullity, so $ \sigma $ maps onto $ \alpha $.
\end{proof}

\section{Problem 5}

\begin{lemma}
    
\end{lemma}

\section{Problem 6}

\begin{lemma}
    $ F $ is a perfect field if and only if every irreducible polynomial $ f $ over $ F $ is separable.
\end{lemma}
\begin{proof}
    This result is trivial if $ char(F) = 0 $. Thus, assume $ char(F) = p > 0 $.
    
    Assume $ F $ is perfect. Let $ f \in F[x] $ be irreducible.

    Assume by contradiction that $ f $ is not separable. Then, there's some $ g(x) \in F[x] $
    such that $ f(x) = g(x^{p}) $ such that $ g(x) $ is irreducible.

    Let $ g(x) = a_{n}x^{n} + \cdots + a_{1}x^{1} + a_{0} $.
    Then, $ g(x^{p}) = a_{n}x^{np} + \cdots + a_{1}x^{p} + a_{0} $.
    Since $ F $ is perfect, for every $ i \in \{1,\ldots, n \} $, we have some $ b_{i} \in F $
    such that $ b_{i}^{p} = a_{i} $.

    Then, $ g(x^{p}) = b_{n}^{p}x^{np} + \cdots + b_{1}^{p}x^{p} + b_{0}^{p}  = (b_{n}x^{n} + \cdots + b_{1}x + b_{0})^{p} $,
    which contradicts the irreducibility of $ F $.

    Let's now prove the backwards direction. We'll use the contrapositive.
    Assume $ F $ is not perfect. Then, there's some $ x \in F $ such that
    it doesn't have a $ p $th root. Then, $ f(t) = t^{p} - x $ is irreducible
    over $ F[x] $ since we can't factor it as $ (t - x)^{p} $, but this is the factorization
    in the splitting field which is unique. But then, $ f $ is not separable and irreducible.
\end{proof}

The lemma above implies that every algebraic extension of a perfect field is perfect since
every irreducible polynomial over in $ K[x] $ is separable over $ K $ if and only
if it is separable over $ F $.

\section{Problem 7}

Let $ K/F $ be a finite field extension.

Part (a) follows immediately by the additivity of trace for matrices and multiplicativity of
the determinant.

\section{Problem 9}



\end{document}
